% VI48-DETAILED-HINTS
% VI48-DETAILED-HINT-01
\paragraph{Exercise VI.48.1.}
\begin{hint}
% PEDAGOGY-ENRICHED-VI
Use the sheaf kernel and injectivity on sections.
\end{hint}

% VI48-DETAILED-HINT-02
\paragraph{Exercise VI.48.2.}
\begin{hint}
% PEDAGOGY-ENRICHED-VI
Use a quotient sheaf whose global lifting obstruction is nonzero.
\end{hint}

% VI48-DETAILED-HINT-03
\paragraph{Exercise VI.48.3.}
\begin{hint}
% PEDAGOGY-ENRICHED-VI
Expand the three differences on a triple overlap.
\end{hint}

% VI48-DETAILED-HINT-04
\paragraph{Exercise VI.48.4.}
\begin{hint}
% PEDAGOGY-ENRICHED-VI
Write the new lift as \(t_i+b_i\) with \(b_i\) kernel-valued.
\end{hint}

% VI48-DETAILED-HINT-05
\paragraph{Exercise VI.48.5.}
\begin{hint}
% PEDAGOGY-ENRICHED-VI
Restrict one global lift to every member of the cover.
\end{hint}

% VI48-DETAILED-HINT-06
\paragraph{Exercise VI.48.6.}
\begin{hint}
% PEDAGOGY-ENRICHED-VI
Correct local lifts by a kernel-valued \(0\)-cochain.
\end{hint}

% VI48-DETAILED-HINT-07
\paragraph{Exercise VI.48.7.}
\begin{hint}
% PEDAGOGY-ENRICHED-VI
Start in degree zero and include the first connecting map.
\end{hint}

% VI48-DETAILED-HINT-08
\paragraph{Exercise VI.48.8.}
\begin{hint}
% PEDAGOGY-ENRICHED-VI
Use exactness at \(H^0(\mathcal F'')\).
\end{hint}

% VI48-DETAILED-HINT-09
\paragraph{Exercise VI.48.9.}
\begin{hint}
% PEDAGOGY-ENRICHED-VI
Identify the middle sheaf with a direct sum.
\end{hint}

% VI48-DETAILED-HINT-10
\paragraph{Exercise VI.48.10.}
\begin{hint}
% PEDAGOGY-ENRICHED-VI
A sheaf epimorphism supplies local, not fixed-open, lifts.
\end{hint}

% VI48-DETAILED-HINT-11
\paragraph{Exercise VI.48.11.}
\begin{hint}
% PEDAGOGY-ENRICHED-VI
Lift a cocycle, differentiate, and descend to the kernel.
\end{hint}

% VI48-DETAILED-HINT-12
\paragraph{Exercise VI.48.12.}
\begin{hint}
% PEDAGOGY-ENRICHED-VI
Push local lifts through the morphism of short exact sequences.
\end{hint}

% VI48-DETAILED-HINT-13
\paragraph{Exercise VI.48.13.}
\begin{hint}
% PEDAGOGY-ENRICHED-VI
Use the ideal sequence and \(\mathcal I_D\simeq\mathcal O_X(-D)\).
\end{hint}

% VI48-DETAILED-HINT-14
\paragraph{Exercise VI.48.14.}
\begin{hint}
% PEDAGOGY-ENRICHED-VI
Line bundles are locally free and therefore flat.
\end{hint}

% VI48-DETAILED-HINT-15
\paragraph{Exercise VI.48.15.}
\begin{hint}
% PEDAGOGY-ENRICHED-VI
Read the connecting map out of the divisor long exact sequence.
\end{hint}

% VI48-DETAILED-HINT-16
\paragraph{Exercise VI.48.16.}
\begin{hint}
% PEDAGOGY-ENRICHED-VI
Tensor the point ideal sequence by \(\mathcal O(n)\).
\end{hint}

% VI48-DETAILED-HINT-17
\paragraph{Exercise VI.48.17.}
\begin{hint}
% PEDAGOGY-ENRICHED-VI
Recall that skyscraper sheaves are flabby.
\end{hint}

% VI48-DETAILED-HINT-18
\paragraph{Exercise VI.48.18.}
\begin{hint}
% PEDAGOGY-ENRICHED-VI
Use surjectivity of evaluation for nonnegative twists.
\end{hint}

% VI48-DETAILED-HINT-19
\paragraph{Exercise VI.48.19.}
\begin{hint}
% PEDAGOGY-ENRICHED-VI
For negative twists the two \(H^0\)-groups vanish.
\end{hint}

% VI48-DETAILED-HINT-20
\paragraph{Exercise VI.48.20.}
\begin{hint}
% PEDAGOGY-ENRICHED-VI
Insert acyclicity of the flabby middle sheaf into the long exact sequence.
\end{hint}

% VI48-DETAILED-HINT-21
\paragraph{Exercise VI.48.21.}
\begin{hint}
% PEDAGOGY-ENRICHED-VI
Read the flabby embedding exact sequence in degree zero.
\end{hint}

% VI48-DETAILED-HINT-22
\paragraph{Exercise VI.48.22.}
\begin{hint}
% PEDAGOGY-ENRICHED-VI
Use the alternating dimension sum of a finite exact sequence.
\end{hint}

% VI48-DETAILED-HINT-23
\paragraph{Exercise VI.48.23.}
\begin{hint}
% PEDAGOGY-ENRICHED-VI
Distinguish which two of the three sheaves are assumed acyclic.
\end{hint}

% VI48-DETAILED-HINT-24
\paragraph{Exercise VI.48.24.}
\begin{hint}
% PEDAGOGY-ENRICHED-VI
Vanishing collapses neighboring terms of the long exact sequence.
\end{hint}
